\section{Discusiones}\label{sec:discusiones}

Los resultados experimentales reportados en la Sección~\ref{sec:resultados} permiten evaluar de manera integral el desempeño físico de la plataforma mecatrónica frente a las predicciones analíticas del modelo de Euler-Lagrange, la simulación numérica no lineal y los cuatro objetivos específicos formulados en la Sección~\ref{sec:introduccion}. A continuación se examinan las razones físicas y matemáticas que explican las coincidencias paramétricas y las discrepancias dinámicas observadas entre el diseño asistido por computador (CAD), la simulación a \qty{1.0}{\kilo\hertz} y el prototipo real, así como el diagnóstico de estabilidad en frecuencia y su posicionamiento frente a la literatura científica.

\subsection{Interpretación Comparativa: Teórico/CAD, Simulación y Hardware}\label{subsec:comparacion_general}

La Tabla~\ref{tab:comparacion_general} consolida la comparación cuantitativa entre los parámetros nominales extraídos del ensamble sólido en FreeCAD~1.1, las métricas del entorno de simulación no lineal ejecutado a \qty{1.0}{\kilo\hertz} y las mediciones físicas adquiridas en tiempo real sobre el banco experimental a través del enlace serial \texttt{COM13}.

\begin{table}[!t]
\centering
\caption{Comparación consolidada entre el modelo teórico/CAD, la simulación no lineal a \qty{1.0}{\kilo\hertz} y el hardware real (\texttt{COM13}).}
\label{tab:comparacion_general}
\resizebox{\columnwidth}{!}{%
\begin{tabular}{@{}lcccc@{}}
\toprule
\textbf{Parámetro / Métrica} & \textbf{Teórico / CAD} & \textbf{Simulación} & \textbf{Hardware Real} & \textbf{Error Rel.} \\
\midrule
Ganancia $K_v$ [\unit{(\radian\per\second)\per\volt}] & $3.40$ & $3.372$ & $3.372$ & $-0.8\%$ \\
Constante $\tau_m$ [\unit{\second}] & $0.1200$ & $0.1175$ & $0.1175$ & $-2.1\%$ \\
Frecuencia $\omega_n$ [\unit{\radian\per\second}] & $8.95$ & $8.8484$ & $8.8484$ & $-1.1\%$ \\
Amortiguamiento $\zeta$ [--] & $0.0500$ & $0.0649$ & $0.0649$ & $+29.8\%$ \\
Factor $K_{\text{acople}}$ [--] & $0.850$ & $0.8737$ & $0.8737$ & $+2.8\%$ \\
Inercia $J_{1,\text{piv}}$ [\unit{\kilogram\meter\squared}] & \num{2.100e-4} & \num{2.044e-4} & \num{2.044e-4} & $-2.7\%$ \\
Umbral $u_{\text{breakward}}^{-}$ [\unit{\volt}] & $-1.20$ & $-2.52$ & $-2.52$ & $-110.0\%$ \\
Tiempo $t_{\text{swing}}$ [\unit{\second}] & $\le 1.50$ & $1.48$ & $0.98$ & $-33.8\%$ \\
$\mathrm{RMS}(\tilde{\alpha})$ [\unit{\degree}] & --- & $1.15$ & $1.99$ & $+73.0\%$ \\
$\mathrm{RMS}(\theta)$ [\unit{\degree}] & --- & $4.85$ & $9.31$ & $+92.0\%$ \\
$\mathrm{RMS}(u)$ [\unit{\volt}] & --- & $0.58$ & $0.73$ & $+25.9\%$ \\
Caídas en equilibrio & $0$ & $0$ & $0$ & $0.0\%$ \\
\bottomrule
\end{tabular}%
}
\end{table}

El análisis de la Tabla~\ref{tab:comparacion_general} revela tres regímenes físicos claramente diferenciados. En primer lugar, los parámetros gobernados por la geometría macroscópica, la distribución de masa y la conversión electromecánica ($K_v$, $\tau_m$, $\omega_n$, $K_{\text{acople}}$ y $J_{1,\text{piv}}$) exhiben una discrepancia relativa inferior al \qty{3}{\percent} (con un máximo de $+2.8\%$ en el factor de acoplamiento inercial cruzado) entre el cálculo tensorial de FreeCAD~1.1 definido por~\eqref{eq:inercias_equiv} y la identificación física sobre las $3\,726$ muestras experimentales. Esta alta fidelidad geométrica e inercial valida directamente la estrategia de diseño para manufactura aditiva (DfAM) descrita en la Sección~\ref{sec:metodologia} (Subsección~\ref{subsec:dfam}): la eliminación total de material de soporte (índice \texttt{SUP} nulo) y la ausencia de interferencias volumétricas (cero colisiones entre los $23$ cuerpos del ensamble) evitaron sobrepesos parásitos, delaminaciones y deformaciones térmicas en las piezas impresas en PLA, garantizando que los tensores de inercia físicos coincidieran con el gemelo geométrico.

En segundo lugar, las propiedades disipativas y de umbral de conducción presentan desviaciones significativas respecto a las estimaciones teóricas ideales. El factor de amortiguamiento viscoso adimensional $\zeta$ experimentó un incremento del $+29.8\%$ respecto al valor nominal de diseño ($0.0649$ frente a $0.0500$). Físicamente, este amortiguamiento suplementario proviene de la película lubricante y los sellos laterales del rodamiento rígido de bolas \texttt{608ZZ} en el pivote secundario, sumados al arrastre aerodinámico del contrapeso cilíndrico bipartito M3 durante las oscilaciones de gran amplitud. De forma análoga, el umbral negativo de ruptura de movimiento $u_{\text{breakward}}^{-}$ registró una desviación del $-110.0\%$ (\qty{-2.52}{\volt} reales frente a \qty{-1.20}{\volt} nominales). Este fenómeno obedece a la asimetría constructiva entre las etapas superior (PNP) e inferior (NPN) de los transistores bipolares en configuración Darlington dentro del puente~H \texttt{L298N}~\cite{L298N_Datasheet}, combinada con la fricción estática direccional de las escobillas de grafito del motor DC y el engrane inicial de la primera etapa planetaria.

En tercer lugar, al contrastar las trayectorias en lazo cerrado (Fig.~\ref{fig:sim_vs_exp}), se observa que la maniobra física de levantamiento (\textit{swing-up}) alcanzó la ventana de captura en \qty{0.980}{\second}, resultando un \qty{33.8}{\percent} más rápida que la predicción del modelo simulado (\qty{1.48}{\second}). La explicación física de esta aceleración radica en la interacción entre el estado de reposo inicial con fricción seca de Coulomb~\eqref{eq:friccion_mixta} y el pulso determinista de arranque: mientras que el integrador numérico continuo debe desarrollar un primer semiciclo de muy baja amplitud para romper la simetría energética desde el equilibrio inferior, en el prototipo real la liberación súbita de la esticción al superar la rampa de zona muerta~\eqref{eq:comp_l298n} inyecta un impulso angular neto que omite dicho semiciclo preliminar y engancha de inmediato el bombeo resonante en cuadratura. Por el contrario, una vez en la fase de regulación vertical (\texttt{BALANCE}), los valores eficaces experimentales $\mathrm{RMS}(\tilde{\alpha})$ de \qty{1.99}{\degree}, $\mathrm{RMS}(\theta)$ de \qty{9.31}{\degree} y $\mathrm{RMS}(u)$ de \qty{0.73}{\volt} superan a sus contrapartes simuladas (\qty{1.15}{\degree}, \qty{4.85}{\degree} y \qty{0.58}{\volt}, respectivamente). Este ensanchamiento de la banda de operación estacionaria es la manifestación física directa de tres perturbaciones estocásticas y estructurales del banco real: la zona libre de inversión (\textit{backlash}) del tren epicicloidal $80:1$, el piso de cuantización de 12~bits de los sensores magnéticos AS5600~\cite{ams2018} y el par restaurador no lineal ejercido por el arnés de cuatro conductores del bus $\text{I}^2\text{C}_2$. No obstante, ambas realizaciones mantienen exactamente cero caídas durante toda la ventana de prueba, confirmando la robustez de la ley implementada.

\subsection{Diagnóstico de Inestabilidad Offline (L-BFGS-B) frente a M4 Ultra-Silk}\label{subsec:diagnostico_lbfgsb}

Un hallazgo fundamental del proceso de desarrollo fue la divergencia experimental observada al evaluar en el hardware la parametrización preliminar obtenida mediante optimización numérica fuera de línea con el algoritmo \texttt{L-BFGS-B} (con ganancias de péndulo $K_{p,\alpha}$ de $52.01$, $K_{i,\alpha}$ de $16.06$ y $K_{d,\alpha}$ de $2.23$, junto con ganancias de brazo $K_{p,\theta}$ de $-2.73$, $K_{i,\theta}$ de $-2.00$ y $K_{d,\theta}$ de $-2.58$). Aunque dicho vector de ganancias minimizaba un funcional cuadrático sobre el modelo ideal sin retardo ni holgura, su implantación en el microcontrolador STM32F446RE provocó la pérdida inmediata de estabilidad mediante una oscilación creciente cercana a \qty{5.24}{\hertz}. La sustitución de este conjunto preliminar en la planta discreta identificada~\eqref{eq:zoh_discreto} demuestra que el lazo cerrado desplazaba un par de polos complejos conjugados hacia el semiplano derecho continuo, con parte real positiva de \qty{+6.48}{\radian\per\second} y frecuencia amortiguada de \qty{32.91}{\radian\per\second}, equivalentes en el plano discreto a un módulo de $1.0065$ fuera del círculo unitario. El análisis físico-matemático permite aislar tres mecanismos causales de esta inestabilidad:

\begin{enumerate}
    \item \textbf{Competencia integral en arquitectura SIMO:} El sistema de Furuta posee una única entrada de control $u(t)$ para regular simultáneamente dos grados de libertad ($\tilde{\alpha}$ y $\theta$). Como lo demuestra la función de transferencia linealizada~\eqref{eq:ft_pendulo_up}, la planta presenta un cero de transmisión en el origen que bloquea el paso de componente continua desde la tensión de armadura hacia la aceleración angular de la barra en régimen estacionario. En presencia de una imperfección geométrica de verticalidad o un sesgo de calibración de apenas \qty{0.2}{\degree}, la activación simultánea de una ganancia integral de péndulo $K_{i,\alpha}$ positiva ($16.06$) y una ganancia integral de brazo $K_{i,\theta}$ no nula ($-2.00$) establece dos acumuladores integrales acoplados que compiten por el mismo actuador. Esta pugna interna induce un ciclo de saturación cruzada y reduce drásticamente el margen de fase requerido para estabilizar el polo gravitacional inestable $s_4$ de \qty{+8.293}{\radian\per\second} caracterizado en~\eqref{eq:polos_pendulo_simb}.
    \item \textbf{Exceso de ganancia proporcional y derivativa frente a retardos de fase no modelados:} De acuerdo con las ecuaciones dinámicas acopladas~\eqref{eq:pendulo_linealizado} y~\eqref{eq:edo_brazo}, una ganancia proporcional $K_{p,\alpha}$ de $52.01$ (equivalente a $3.6$ veces el valor final estable de $14.50$) combinada con una magnitud derivativa de brazo de $2.58$ ($3.0$ veces superior a la magnitud final de $0.85$) exige variaciones instantáneas de par que superan el ancho de banda mecánico del reductor. Al interactuar con el retardo de fase introducido por la holgura del tren de engranajes y con la amplificación de alta frecuencia de la diferenciación discreta~\cite{Merry2010}, el incremento excesivo de ganancia desplaza las ramas del lugar geométrico de las raíces hacia el exterior del círculo unitario en la frecuencia crítica de \qty{5.24}{\hertz}.
    \item \textbf{Ausencia de filtrado estructural selectivo y de limitación de pendiente:} El esquema optimizado fuera de línea carecía de atenuación específica sobre la resonancia mecánica de \qty{11.0}{\hertz} y de restricción en la derivada temporal del mando. En consecuencia, cada inversión de signo del error excitaba impactos diente contra diente en la reductora $80:1$, saturando de forma alternante el puente \texttt{L298N} y realimentando el ciclo límite divergente.
\end{enumerate}

La arquitectura definitiva \texttt{M4 Ultra-Silk}~\eqref{eq:m4_dual} resuelve simultáneamente estas tres patologías estructurales. En primer término, anula estructuralmente la acción integral sobre la coordenada pendular ($K_{i,\alpha}$ fijada en cero), reservando el canal integral exclusivamente para el lazo lento del brazo horizontal ($K_{i,\theta}$ de $-0.75$, acotado en $\pm\qty{1.6}{\volt}$) a fin de absorber asimetrías de cableado sin disputar la autoridad de fase del estabilizador rápido $\text{PD}_\alpha$. En segundo término, la inserción del filtro pasa-bajos IIR en cascada~\eqref{eq:lpf_iir} con función de transferencia discreta~\eqref{eq:lpf_num_z}, seguida en serie por el filtro biquad \textit{Dual-Notch}~\eqref{eq:dual_notch} conformado por las etapas~\eqref{eq:hz_n11} y~\eqref{eq:hz_n6} bajo la ley de ponderación adaptativa~\eqref{eq:mezcla_dist}, suprime selectivamente la energía en las bandas de \qty{11.0}{\hertz} y \qty{6.2}{\hertz} sin degradar el margen de fase en la banda de cruce del péndulo. Finalmente, la moderación de las ganancias proporcional y derivativa bajo saturación subordinada del brazo ($\pm\qty{2.2}{\volt}$) junto con el limitador adaptativo de pendiente reubica la totalidad de los cinco polos dominantes de lazo cerrado en el interior del círculo unitario, tal como se verificó analíticamente en~\eqref{eq:polos_m4}, alcanzando un radio espectral $\rho(z)$ de $0.99934$ estrictamente menor que la unidad~\cite{Ogata2010,Astrom1997,Franklin2015}.

\subsection{Fuentes Físicas de No Idealidad, Contraste con la Literatura y Alcance de Objetivos}\label{subsec:fuentes_no_idealidad}

La operación continua del prototipo a \qty{1.0}{\kilo\hertz} permitió aislar y cuantificar seis fuentes físicas de no idealidad inherentes a los sistemas mecatrónicos subactuados construidos con componentes accesibles, así como validar la eficacia de las contramedidas algorítmicas adoptadas:

\begin{enumerate}
    \item \textbf{Holgura mecánica (\textit{backlash}) del reductor epicicloidal $80:1$:} Durante las inversiones de par alrededor del equilibrio vertical, el juego libre entre los dientes planetarios desacopla momentáneamente al motor del brazo horizontal, generando un ciclo límite torsional a \qty{11.0}{\hertz} acompañado de una oscilación secundaria a \qty{6.2}{\hertz}. La cascada \textit{Dual-Notch}~\eqref{eq:hz_n11} y~\eqref{eq:hz_n6} con ganancia de mezcla adaptativa~\eqref{eq:mezcla_dist} suprime el modo primario de \qty{11.0}{\hertz} y atenúa en \qty{-6.0}{\decibel} el modo secundario de \qty{6.2}{\hertz} en quietud vertical, extinguiendo la vibración sin inducir el retardo de fase que produciría un filtro pasa-bajos convencional de baja frecuencia de corte~\cite{Oppenheim2010}.
    \item \textbf{Asimetría de conducción BJT Darlington en el puente L298N:} A diferencia de los puentes integrados basados en transistores MOSFET de baja resistencia de conducción, la etapa bipolar del \texttt{L298N}~\cite{L298N_Datasheet} introduce una caída de saturación colector-emisor $V_{\text{sat}}$ cercana a \qty{2.0}{\volt} y una zona muerta direccional severamente asimétrica ($+10.0\%$ frente a $-21.0\%$ de ciclo útil PWM). La compensación mediante rampa continua suave~\eqref{eq:comp_l298n} con banda silenciosa de \qty{0.06}{\volt} y rodilla $V_k=\qty{0.65}{\volt}$ restituye la autoridad de control en la vecindad del origen sin desencadenar el castañeteo (\textit{chattering}) de conmutación característico de las funciones signo discontinuas en moduladores PWM discretos~\cite{Petric2022}.
    \item \textbf{Ruido de cuantización derivativa en codificadores magnéticos de 12~bits:} Con una resolución angular de $4\,096$ cuentas por revolución (\qty{1.534e-3}{\radian} o \qty{0.0879}{\degree} por bit menos significativo) y un periodo de muestreo isócrono $T_s$ de \qty{1.0}{\milli\second}, la evaluación numérica de la cota teórica de diferenciación hacia atrás~\eqref{eq:diff_ruido_teoria} arroja impulsos espurios de velocidad de magnitud~\cite{Merry2010}:
    \begin{equation}
    \Delta\dot{\alpha}_{\text{ruido}} = \pm \frac{\qty{1.534e-3}{\radian}}{\qty{1.0e-3}{\second}} = \pm\qty{1.534}{\radian\per\second}
    \label{eq:ruido_derivativo}
    \end{equation}
    De haberse inyectado directamente sobre la ganancia derivativa del péndulo, las espigas descritas por~\eqref{eq:ruido_derivativo} habrían inducido conmutaciones parásitas superiores a \qty{2.3}{\volt} en cada milisegundo. Su mitigación efectiva se logró mediante dos vías complementarias: en el eje pendular, la doble cascada IIR~\eqref{eq:lpf_num_z} impone una caída asintótica de \qty{-40}{\decibel} por década por encima de $f_c\approx\qty{15.8}{\hertz}$~\cite{Oppenheim2010}; en el eje del brazo, el montaje mecánico del primer sensor AS5600 directamente sobre el rotor de alta velocidad multiplica por $80$ la densidad espacial de muestreo tras la reducción~\cite{ams2018}, reduciendo el escalón de cuantización cinemática a una fracción despreciable (\qty{0.0011}{\degree}).
    \item \textbf{Fricción mixta viscosa y seca de Coulomb en el pivote \texttt{608ZZ}:} La identificación experimental del modelo disipativo~\eqref{eq:friccion_mixta} reveló un término de fricción seca $c_c$ de \qty{0.6179}{\radian\per\second\squared} en la articulación libre del péndulo. Dado que este par resistivo se opone de forma discontinua a cualquier inversión de velocidad angular $\dot{\alpha}$, establece un umbral físico irreductible debajo del cual la barra se adhiere momentáneamente antes de ser corregida por el par inercial del brazo, explicando por qué el error eficaz estacionario $\mathrm{RMS}(\tilde{\alpha})$ converge a \qty{1.99}{\degree} en el hardware frente a \qty{1.15}{\degree} en la simulación.
    \item \textbf{Par elástico residual del arnés distal $\text{I}^2\text{C}_2$:} Los cuatro conductores flexibles que alimentan y comunican al sensor AS5600 del pivote secundario introducen una rigidez torsional dependiente de la posición azimutal $\theta$. El integrador acotado del controlador $\text{PID}_\theta$~\eqref{eq:m4_dual} absorbe este par restaurador como una perturbación exógena cuasiestacionaria, manteniendo al brazo confinado alrededor de su referencia sin desestabilizar la barra vertical.
    \item \textbf{Desacoplamiento inercial reflejado por la relación de reducción $N_g$:} A partir de las inercias equivalentes deducidas en~\eqref{eq:inercias_equiv}, el factor cuadrático $N_g^2$ de $6\,400$ magnifica la inercia rotórica reflejada $J_{\text{ref}}$ en varios órdenes de magnitud por encima del término de acoplamiento inercial cruzado entre el brazo y el péndulo, satisfaciendo la razón adimensional:
    \begin{equation}
    \frac{J_{\text{ref}}}{m_1 L_0 l_1} \approx 103.4 \gg 1
    \label{eq:desacople_ratio}
    \end{equation}
    La desigualdad constructiva expresada en~\eqref{eq:desacople_ratio} demuestra físicamente que el par de reacción gravitacional y centrífugo ejercido por el eslabón pendular representa menos del \qty{1}{\percent} de la inercia efectiva vista por el eje de salida de la reductora. Esta propiedad estructural explica por qué la dinámica de velocidad del actuador pudo identificarse con un coeficiente de determinación $R^2$ del $98.51\%$ mediante la función de transferencia lineal de primer orden~\eqref{eq:gv_identificada}.
\end{enumerate}

Al contrastar estos hallazgos con la literatura clásica sobre péndulos rotacionales subactuados, emergen aportaciones teóricas y prácticas relevantes. Los esquemas pioneros de levantamiento por realimentación de pseudoestado de Furuta, Yamakita y Kobayashi~\cite{furuta1991}, de linealización parcial por bombeo energético de Spong~\cite{spong1995} y de control de energía por Lyapunov de {\AA}str{\"o}m y Furuta~\cite{astrom2000} fueron formulados bajo la premisa de una rotación azimutal libre de restricciones mecánicas, apoyada habitualmente en anillos rozantes industriales de bajo par. En plataformas académicas sin anillos colectores, la aplicación directa de tales leyes provoca el enrollamiento destructivo del cableado del sensor distal durante la fase de aceleración. Los resultados de la Fig.~\ref{fig:exp_validacion} demuestran que la ley de bombeo resonante en cuadratura implementada, dotada de seguimiento proporcional-derivativo con saturación escalonada y carrera acotada, logra completar el levantamiento autónomo en \qty{0.980}{\second} restringiendo simultáneamente la excursión máxima comandada del brazo horizontal a $\pm\qty{50}{\degree}$.

Asimismo, en relación con el modelado analítico no lineal de Cazzolato y Prime~\cite{cazzolato2011}, el presente trabajo confirma que la formulación completa de Euler-Lagrange conserva una capacidad predictiva sobresaliente incluso en estructuras modulares manufacturadas íntegramente en PLA mediante impresión 3D FDM, siempre que el diseño mecánico elimine holguras estructurales y soportes adheridos (índice \texttt{SUP} nulo) y que las no linealidades de fricción mixta~\eqref{eq:friccion_mixta} y zona muerta del actuador~\eqref{eq:comp_l298n} se incorporen empíricamente al modelo.

En el ámbito del control digital con saturación y conmutación de modos, la integración de las técnicas de condicionamiento y transferencia sin sacudidas (\textit{bumpless transfer}) de {\AA}str{\"o}m y Rundqwist~\cite{AstromAntiwindup1989} y de Peng, Vran{\v{c}}i{\'c} y Hanus~\cite{Peng1996} permitió ejecutar la transición entre la ley energética no lineal (\texttt{SWING\_UP}) y el estabilizador lineal (\texttt{BALANCE\_PID}) sin discontinuidades de par que abatieran el eslabón al cruzar el cono de captura de $\pm\qty{20}{\degree}$. Adicionalmente, la extensión discreta del integrador no lineal de Clegg~\cite{clegg1958} mediante la regla de reinicio direccional con histéresis~\eqref{eq:clegg_pos} y freno electromagnético activo resolvió el problema del sobreimpulso por memoria integral durante el retorno al origen (\texttt{STOPHOME}), comprimiendo el sobreimpulso del ensayo servo \texttt{CHAR5B} al $0.29\%$ y alcanzando una exactitud final de parqueo de \qty{-0.25}{\degree}.

Desde una perspectiva práctica, conviene señalar tres limitaciones operativas del sistema desarrollado que delimitan sus líneas de mejora futura:
\begin{itemize}
    \item \textbf{Eficiencia energética y deriva térmica del puente BJT:} La caída bipolar del módulo \texttt{L298N}~\cite{L298N_Datasheet} disipa potencia reactiva en forma de calor y desplaza levemente los umbrales de zona muerta tras sesiones prolongadas de operación continua; su sustitución por un puente MOSFET de conmutación rápida con sensado de corriente en lazo interno permitiría comandar directamente par electromagnético en lugar de tensión de armadura.
    \item \textbf{Piso de resolución en la articulación pendular:} A diferencia del sensor del motor beneficiado por el factor $80:1$, el codificador AS5600 distal opera en acoplamiento directo $1:1$ con un límite físico de \qty{0.0879}{\degree} por cuenta~\cite{ams2018}, lo cual acota la ganancia derivativa máxima admisible antes de amplificar ruido de cuantización~\cite{Merry2010}.
    \item \textbf{Sensibilidad numérica de los filtros biquad de alto factor de calidad:} La ubicación de los polos de las muescas~\eqref{eq:hz_n11} y~\eqref{eq:hz_n6} a una distancia radial inferior a $0.08$ respecto a la circunferencia unitaria exige aritmética de punto flotante de doble precisión (64~bits) en la estructura TDF-II sobre el procesador Cortex-M4F~\cite{STM32F446RE_Datasheet,Oppenheim2010} para prevenir ciclos límite por truncamiento de coeficientes.
\end{itemize}

En síntesis, la evidencia experimental discutida verifica el cumplimiento cabal de los cuatro objetivos específicos trazados en la Sección~\ref{sec:introduccion}: se concretó una arquitectura mecatrónica fabricable sin soportes con alta correspondencia CAD-hardware (Objetivo~1); se identificaron experimentalmente con rigor estadístico las dinámicas lineales y no lineales del conjunto actuador-péndulo (Objetivo~2); se sintetizó e implementó a \qty{1.0}{\kilo\hertz} la arquitectura \texttt{M4 Ultra-Silk} con filtrado \textit{Dual-Notch} adaptativo y servo de Clegg garantizando estabilidad asintótica (Objetivo~3); y se validó el desempeño integral junto al gemelo digital 3D en tiempo real logrando levantamiento autónomo sub-segundo y regulación continua sin caídas (Objetivo~4).
